\documentclass[11pt,a4paper]{report}
\usepackage[T1]{fontenc}
\usepackage[utf8]{inputenc}
\usepackage{lmodern}
\usepackage{geometry}
\usepackage{setspace}
\usepackage{graphicx}
\usepackage{amsmath,amssymb}
\usepackage{hyperref}

\geometry{margin=1in}
\onehalfspacing

\title{KDEllipspy's Documentation}
\date{\today}

\begin{document}

\maketitle
\tableofcontents
\clearpage
\section*{Introduction}
\textit{\textbf{KDEllipspy}} is a Python-based software developed for the kinematic and dynamic inversion of earthquakes. By implementing a physical model that assumes an elliptical slip distribution for the seismic source, the package facilitates the inversion of seismic data to robustly estimate key source parameters.\par
To achieve computational efficiency and physical accuracy, the software leverages two established numerical engines: it utilizes \textbf{AXITRA} to compute Green's functions and convolve the sources for the kinematic modeling, and \textbf{FD3D\_TSN} to handle the rupture and dynamic wave propagation.\par
This documentation provides a comprehensive guide to the software, structured as follows: the theoretical background underlying the models, a detailed user manual, practical usage examples, and technical notes for developers.
\chapter{Theoretical Background}

\section{Problem Definition}
KDEllipspy addresses earthquake source inversion by combining physical modeling and waveform fitting.
The goal is to estimate source parameters that explain observed seismic records.

\section{Forward Modeling Concepts}
The forward problem maps source parameters to synthetic seismograms.
In this framework, source geometry, rupture timing, and source amplitudes are translated into predicted waveforms at receiver stations.

\section{Inversion Perspective}
The inverse problem searches for parameter sets that minimize a misfit between observed and synthetic waveforms.
Depending on the workflow, this can be done with stochastic optimization (e.g., Neighbourhood Algorithm) or Bayesian sampling (e.g., MCMC).

\section{Assumptions and Limitations}
This section should document key assumptions, such as fixed mesh geometry, velocity model choices, and data-window selection for misfit computation.

\section{Kinematic Source Framework}
The kinematic module follows the finite-fault paradigm: the fault plane is discretized into subfaults, and each subfault contributes to the final seismogram through a delayed source history and a propagation operator.
This formulation is rooted in classical source theory and rupture kinematics (Kostrov, 1964; Haskell, 1964; Aki and Richards, 2002).

\subsection{Subfault Parameterization}
Let the fault be discretized into $N_s$ subfaults.
In KDEllipspy, the slip field is controlled by an elliptical parameterization with semi-axes $(a_1,a_2)$, orientation $\theta$, position parameters, maximum slip $d_{\max}$, and rupture velocity $v_r$.
For a subfault center $(x,y)$ in local fault coordinates,
\begin{equation}
d = \left(\frac{x'}{a_1}\right)^2 + \left(\frac{y'}{a_2}\right)^2,
\end{equation}
where $(x',y')$ are coordinates in the rotated ellipse frame.
The local slip factor $f(d)$ is imposed by the chosen shape law (constant, Gaussian-like, or elliptical taper), and the effective slip is
\begin{equation}
D_i = d_{\max} \, f(d_i).
\end{equation}

\subsection{Rupture Timing}
Rupture delay is modeled from hypocentral distance and rupture speed:
\begin{equation}
t_{r,i} = \frac{\|\mathbf{x}_i - \mathbf{x}_h\|}{v_r},
\end{equation}
with $\mathbf{x}_h$ the hypocenter and $\mathbf{x}_i$ the subfault location.
This defines the source activation chronology over the mesh.

\subsection{Moment and Waveform Synthesis}
For each subfault, scalar moment is proportional to rigidity, area, and slip:
\begin{equation}
M_{0,i} = \mu_i A_i D_i.
\end{equation}
Observed components are synthesized through linear superposition of Green's functions and source histories:
\begin{equation}
u_r(t) = \sum_{i=1}^{N_s}\sum_{k=1}^{N_b}\left(G_{r,i,k} * q_{i,k}\right)(t),
\end{equation}
where $r$ is receiver/component, $k$ denotes source basis terms, and $*$ is temporal convolution.
This representation matches the standard linear elastodynamic framework used in reflectivity/discrete-wavenumber approaches (e.g., Bouchon, 1981; Cotton and Coutant, 1997).

\subsection{Misfit in Inversion}
Model ranking is based on normalized waveform misfit between observed and synthetic data, using P/S windows and rotated components.
This follows the classical finite-fault inversion logic in which source parameters are constrained by waveform similarity in selected time windows (e.g., Hartzell and Heaton, 1983; Kikuchi and Kanamori, 1991).

\chapter{User Manual}

\section{Installation}
Describe installation requirements (Python version, dependencies, and optional tools), and include project setup commands.

\section{Project Structure}
Explain the repository layout, including \texttt{kdellipspy/}, \texttt{Dynamic\_inversion/}, and documentation folders.

\section{Input Data}
Document required input files, expected formats, units, and naming conventions.

\section{Running an Inversion}
Provide the command-line workflow for running kinematic and dynamic workflows, including required flags and output directories.

\section{Outputs}
Describe generated artifacts such as JSON/CSV results, figures, and intermediate files.


\chapter{Usage Examples}

\section{Example 1: Basic Kinematic Inversion}
Provide a minimal end-to-end run with default parameters and expected outputs.

\section{Example 2: Tuning Inversion Parameters}
Show how to modify search settings (iterations, proposal scales, or sampling choices) and discuss practical effects.

\section{Example 3: Dynamic Legacy Workflow}
Document how to execute the legacy dynamic workflow and where to inspect outputs.

\section{Interpreting Results}
Summarize how to read misfit evolution, best-model parameters, and waveform fit plots.


\chapter{Developer Notes}

\section{Software Architecture}
Explain the separation between core modules, inversion modules, input parsing, and workflows.

\section{Performance Considerations}
Document known bottlenecks and optimization strategies, such as Green's function caching and source indexing consistency.

\section{Validation Strategy}
Describe how to validate physics and implementation changes using synthetic tests and reproducible runs.

\section{Roadmap}
List planned features and refactoring goals (for example, dynamic integration and unified inversion interfaces).



\end{document}
